\section{Introduction}
\label{sec:intro}

Self-evolving agents and reusable ``Skills'' are increasingly marketed as engines of
discovery. Systems such as Voyager~\citep{voyager}, Promptbreeder~\citep{promptbreeder}, the
Darwin G\"odel Machine~\citep{dgm}, and the SKILL.md agent-skills abstraction accumulate
procedural knowledge or rewrite their own prompts, and the self-evolving-agents
survey~\citep{selfevolvesurvey} frames this as a paradigm shift from offline pretraining
toward online, lifelong self-modification. Yet almost all of these methods are
\emph{black-box wrappers around a frozen base model}: they add no weights and no training
signal beyond what a verifier or environment already provides. This raises a sharp empirical
question that determines whether ``self-improving AI'' is a path to new science or a plateau
at the base model's latent ceiling.

{\bf When a skill helps, why does it help?} An agent that fails a hard problem and then
succeeds after being given (or self-developing) a skill may have (a) genuinely \emph{acquired}
a new procedure it did not possess --- the ``invented the attention mechanism'' case --- or
(b) merely \emph{recalled} a latent, pretraining-adjacent concept via better prompting. These
are the difference between knowledge creation and prompt optimization, and they call for
opposite conclusions about the field. But standard evaluation cannot separate them. Most
training samples are seen once, so there is no over-fitting fingerprint to detect, and
subjective novelty judgments are actively misleading: HindSight~\citep{hindsight} shows that
LLM and human novelty ratings \emph{anti-correlate} ($\rho=-0.29$) with real downstream impact.

{\bf The gap.} Every existing skill ablation measures only pass-rate \emph{with} versus
\emph{without} a skill. SkillsBench~\citep{skillsbench} shows human-curated skills add
$+16.6$ points while agent self-generated skills fall \emph{below} baseline, and a clean
data-science ablation~\citep{dsskills} finds LLM-generated skills statistically
indistinguishable from a token-matched irrelevant control. These are valuable, but none
isolates \emph{whether the model could have produced the skill itself}, and none has a
substrate where ``novel versus latent'' is ground-truth-known. Without such a substrate,
novelty claims rest on judgments that we already know are untrustworthy.

{\bf Our approach.} We build a controlled substrate where the ground-truth answer to ``is this
skill novel or latent?'' is \emph{known by construction}, and we define the
Elicitation--Novelty Discriminator (\END), a small set of measurable quantities that classify
\emph{why} a skill helped. The key construction is an invented-language glossary (``\zeth''):
its 24 word mappings are generated from a fixed random seed at run time, so the mapping
provably cannot appear in any pretraining corpus. It is the cleanest lab analog of ``a new
theory the model has not memorized,'' yet its answers are exactly checkable by a deterministic
verifier --- no LLM-as-judge. Alongside it we use counterfactual base-$b$
arithmetic~\citep{reasoningreciting} to probe latent-but-unreliable procedures. The \END
asks three questions of any helpful skill: can the base model \emph{self-generate} a working
version (\SG)? Does the skill beat a \emph{length-matched} irrelevant control (\LCG)? And can a
verifier-only \emph{self-evolution} loop reach it without ground-truth answers?

{\bf What we find.} Across two frontier models (\claude, \gpt) and 3{,}000+ deterministic API
evaluations, a skill delivers a large, content-specific, length-controlled gain in exactly one
place: where the required knowledge is genuinely absent. On \zeth the correct skill lifts
accuracy from $0.00$ to $1.00$ on both models ($p<10^{-4}$, Cohen's $h=3.14$), while the
model's own self-generated skill and a length-matched control both stay at $0.00$. A
five-step self-evolution loop never moves \zeth off $0.00$ under verifier-only feedback across
all five iterations, but climbs to $0.85$ the instant it receives ground-truth answers ---
precisely the AlphaEvolve~\citep{alphaevolve} lesson that the model proposes while an external
verifier discovers. Self-generated skills can even \emph{harm}: on base-9 multiplication the
model's own skill dropped accuracy from $0.85$ to $0.25$. Finally, we stress-test the popular
``fine-tune the skill back into the weights'' probe and show it is real but method-dependent:
new knowledge internalizes and generalizes to held-out phrases ($1.00$) only via
example-distillation, while fine-tuning on the skill \emph{document} leaves accuracy at $0.00$.

Our contributions are:
\begin{itemize}[leftmargin=*,itemsep=0pt,topsep=0pt]
    \item We propose a controlled substrate in which ``novel versus latent'' is known by
    construction, built around a randomly generated invented-language glossary whose answers
    are deterministically verifiable, removing the reliance on subjective novelty judgments.
    \item We introduce the Elicitation--Novelty Discriminator (\END): self-generation utility,
    length-controlled gain, and verifier-only self-evolution, which together classify
    \emph{why} a skill helps rather than only \emph{whether} it helps.
    \item We validate the discriminator on two frontier models over 3{,}000+ deterministic
    evaluations, showing a sharp and consistent signature that separates the three regimes and
    mechanizes four prior findings on a single substrate.
    \item We complement the API study with a GPU internalization probe on \qwen that shows the
    ``fine-tune it back'' test cannot by itself classify a skill, since its verdict flips with
    the fine-tuning recipe.
\end{itemize}

We review related work in \secref{sec:related}, describe the substrate and discriminator in
\secref{sec:method}, present results in \secref{sec:results}, discuss implications and
limitations in \secref{sec:discussion}, and conclude in \secref{sec:conclusion}.
